\chapter{Curvature concentration and boundary energy}\label{sec:examples-willmore-weak-limits}
These examples test the hypotheses in Section \ref{sec:willmore-varifold-energy}. They distinguish the area measure, the squared-curvature measure and the singular part of first variation. The weak framework is that of \cite{Sim18}, Chapter 8, Section 2.

\begin{example}[A collapsing sphere loses mass and concentrates curvature energy]\label{ex:collapsing-sphere-willmore-defect}
Let $S_j=\partial B_{r_j}(0)\subset\mathbb R^3$ with $r_j\downarrow0$ and multiplicity one. Then $V_{S_j}\rightharpoonup0$ as varifolds and $\mathcal W(S_j)=4\pi$, but their curvature-energy measures satisfy
\[
 \frac14|\vec H_j|^2\mathcal H^2\llcorner S_j\rightharpoonup4\pi\delta_0.
\]
Thus the weak limit's zero curvature energy does not account for all the energy in the approximating sequence.
\end{example}
\begin{proof}
The sphere has area $4\pi r_j^2$ and $|\vec H_j|=2/r_j$. Its varifold mass therefore tends to zero, bounding every compactly supported test integral by its supremum times that mass. The energy is $4\pi$. For a continuous compactly supported scalar test $\phi$, changing variables $x=r_j\omega$ gives
\[
 \frac14\int_{S_j}\phi(x)|\vec H_j|^2d\mu_j
 =\int_{S^2}\phi(r_j\omega)d\mu_{S^2}\longrightarrow4\pi\phi(0).
\]
Uniform continuity or dominated convergence proves the limit. The zero varifold has zero first variation and zero extended energy. This is strict lower semicontinuity and concentration, not an equality of limit energies.
\end{proof}

\begin{example}[Energy can escape to infinity]\label{ex:willmore-spheres-escape}
The unit spheres centered at $je_1$ converge locally as varifolds to zero while each has area and Willmore energy $4\pi$. Their curvature-energy measures converge vaguely to zero, rather than to a nonzero point mass.
\end{example}
\begin{proof}
Every fixed compact subset of Euclidean space misses those spheres for all sufficiently large $j$. Hence every compactly supported varifold or scalar measure test is eventually zero. Translation does not change either area or mean curvature, giving the constant totals. The loss occurs at infinity and shows why local convergence alone cannot preserve global mass or energy.
\end{proof}

\begin{example}[A disk has zero classical bending energy but infinite full first-variation energy]\label{ex:disk-singular-first-variation-energy}
Let $V$ be the multiplicity-one varifold of a closed planar disk $D_R\subset\mathbb R^3$, considered in all of $\mathbb R^3$. Its smooth interior mean curvature is zero, but $\mathcal E_{\mathbb R^3}(V)=+\infty$. The singular boundary first variation is essential.
\end{example}
\begin{proof}
The divergence theorem gives $\delta V(X)=\int_{\partial D_R}X\cdot\eta\,ds$. Choose smooth compactly supported ambient fields $X_\varepsilon$ that equal the outward radial unit vector on the boundary circle and are supported, on the disk, in an annulus of width at most $\varepsilon$, with uniformly bounded magnitude. Such fields come from smooth radial cutoffs near $r=R$ and a cutoff in the normal coordinate. Then
\[
 \delta V(X_\varepsilon)=2\pi R,\qquad
 b_\varepsilon=\int|X_\varepsilon|^2d\|V\|=O(\varepsilon).
\]
The energy definition tested with $tX_\varepsilon$ gives a supremum over $t$ equal to $(2\pi R)^2/(4b_\varepsilon)$, which tends to infinity. If instead one works in an ambient open set not meeting the boundary circle, its first variation is zero there. The ambient domain is therefore part of the assertion.
\end{proof}

\begin{example}[Real multiplicity scales the energy]\label{ex:willmore-noninteger-sphere}
For the round sphere varifold $V=q|S_R^2|$ with a constant real multiplicity $q>0$,
\[
 \mathcal E_{\mathbb R^3}(V)=4\pi q.
\]
If $0<q<1$, this is below the smooth closed-surface bound $4\pi$. For every noninteger $q$ the varifold is rectifiable but not integral, and cannot come from a closed immersion counted with ordinary parameter multiplicity. For positive integer $q$ it is integral.
\end{example}
\begin{proof}
Multiplying both first variation and weight by $q$ does not change the representing geometric mean curvature. Its norm is $2/R$ and its weighted area is $4\pi qR^2$. Theorem \ref{thm:varifold-energy-dual-representation} therefore gives $4\pi q$. Its density is the real number $q$, so it is nonintegral for $q\notin\mathbb N$. The case $0<q<1$ demonstrates why no unqualified extension of Theorem \ref{thm:li-yau-smooth-multiplicity} to all nonzero varifolds is being asserted.
\end{proof}
